% 1. ит фром бит (объединие физического и информационного)
% 2. ультраметрические метрики окружают нас
% 3. информационная энтропия в замкнутах пространствах
% 4. бесконечность это -1 (тогда 0+0 != 0). 0 с тз физики это сумма противоположных (в силу свойства аддитивности), а для нас это напряжение
% 5. модулярная арифетика (для нас сейчас самая простая операция это сложение, а нужно чтобы было умножение)
% 6. простые числа самые крутые (2 самое крутое)
% 7. Дуплекс: из физической реальности в отображение на наше восприятие мира и наоборот)

\documentclass[areasetadvanced]{scrartcl}
\usepackage[T2A]{fontenc}
\usepackage[english,russian]{babel}

\usepackage[footskip=1cm,left=25mm, right=15mm, top=20mm, bottom=20mm]{geometry}
\usepackage{setspace}
\usepackage{amsmath, amssymb} 
\usepackage{graphicx}
\usepackage{tikz}
\usetikzlibrary{arrows.meta}
\usepackage{float}
\usepackage{dashrule}
\usepackage{fancyhdr} 
\usepackage{hyperref} 
\usepackage{parskip}
\usepackage{textcomp, enumitem}
\usepackage{indentfirst}
\usepackage{graphicx}
\usepackage{algorithm}
\usepackage{algpseudocode}
\usepackage{array} 
\usepackage{geometry}
\usepackage{afterpage}
\usepackage{minted}
\setcounter{secnumdepth}{3} 
\setcounter{tocdepth}{3}    
\usepackage{listings} 

\newcommand{\icon}[1]{\includegraphics[height=1.2em]{\#1}}

\tikzstyle{block} = [rectangle, rounded corners, minimum width=3cm, minimum height=1cm, text centered, draw=black, fill=lightgray]

\setkomafont{sectioning}{\normalfont\bfseries} 
\setkomafont{section}{\normalfont\Large\bfseries}
\setkomafont{subsection}{\normalfont\large\bfseries}
\setkomafont{subsubsection}{\normalfont\large\bfseries}
\setkomafont{paragraph}{\normalfont\large\bfseries} 

\lstset{
  language=Haskell,
  basicstyle=\ttfamily\small,
  keywordstyle=\color{blue}\bfseries,
  stringstyle=\color{red},
  commentstyle=\color{green!70!black},
  numbers=left,
  numberstyle=\tiny,
  stepnumber=1,
  numbersep=10pt,
  showstringspaces=false,
  breaklines=true,
  frame=single
}

\lstdefinelanguage{Lua}{
    keywords={function, end, if, then, else, elseif, for, while, do, repeat, until, break, return, local, and, or, not, true, false, nil},
    keywordstyle=\color{blue}\bfseries,
    stringstyle=\color{red},
    commentstyle=\color{green!70!black},
    morestring=[s]{"}{"},
    morestring=[s]{'}{'},
    morecomment=[l]{--},
    morecomment=[s]{--[[}{]]},
    basicstyle=\ttfamily\small,
    numbers=left,
    numberstyle=\tiny,
    stepnumber=1,
    numbersep=10pt,
    showstringspaces=false,
    breaklines=true,
    frame=single
}

\setlength{\parindent}{1.25cm}
\setcounter{tocdepth}{3}
\begin{document}
    \sloppy
    \begin{center}
        \Huge{\textbf{7 методических принципов}} 
        \\[20cm]
        \large{Салимли Айзек. Группа 5130201/20101.}
    \end{center}
    \newpage
        \section{It from Bit}
            Фраза принадлежит физику Джону Арчибальду Уилеру (John Archibald Wheeler). 
            Она выражает идею, что физическая реальность (it) происходит из информации (bit).
            Связь физического и информационного проявляется, например, при измерении: 
            физический ввод (напряжение, свет) преобразуется в информационный бит, который может вызывать сдвиг регистра.
            В компьютере информация хранится физически (заряды конденсаторов, состояние триггеров), 
            а операции над битами меняют физическое состояние. Компьютер можно рассматривать как \emph{трансформер}, 
            преобразующий входной поток битов в выходной через физические процессы.

            \textbf{Математическая модель:} любое вычисление представимо как функция 
            $f: \{0,1\}^n \to \{0,1\}^m$, реализуемая композицией логических элементов.

        \section{Ультраметрика}
            Ультраметрика — метрика, удовлетворяющая усиленному неравенству треугольника:
            $d(x,z) \le \max(d(x,y), d(y,z))$. 
            Это приводит к тому, что все треугольники равнобедренные с основанием не больше боковых сторон.

            \textbf{Примеры:}
            \begin{itemize}
                \item Пусть узлы $x, y$ имеют иерархические адреса (например, IP). 
                Определим расстояние $d(x,y)$ как глубину их наименьшего общего предка в дереве адресации:
                \[
                d(x,y) = \text{уровень}(\text{LCA}(x,y)).
                \]
                Тогда $d$ — ультраметрика: $\max(d(x,y), d(y,z)) \ge d(x,z)$.
                Число хопов (пересылок) между узлами может быть одинаковым (например, 5), 
                но реальное время передачи (RTT) — разным из-за задержек и загрузки. 
                Ультраметрика игнорирует эти вариации и даёт дискретную меру, зависящую только от иерархии.
                \item Температура воздуха сама по себе не метрика, но скорость молекул (среднеквадратичная) может задавать метрику 
                в фазовом пространстве.
                \item В компьютерных науках: расстояние между двумя файлами в иерархической файловой системе можно определить 
                как глубину их ближайшего общего предка. Если корень — общий предок всех, то это ультраметрика.
            \end{itemize}
        \section{Информационная энтропия}
                В компьютерных науках энтропия — мера неопределённости источника данных. 
                Формула Шеннона для дискретного источника:
                \[
                H = -\sum_{i=1}^{n} p_i \log_2 p_i \quad \text{(бит)}.
                \]
                Энтропия определяет минимальную среднюю длину кода при оптимальном сжатии. 
                В замкнутых системах (например, ограниченный объём памяти) энтропия ограничена числом бит.
                        
                \textbf{Модель:} случайная величина $X$ с распределением $\{p_i\}$; 
                кодирование сообщений длины $N$ требует не менее $N \cdot H(X)$ бит (теорема Шеннона).
                        
        \section{Бесконечность -1}
                В ограниченной разрядной сетке числа представляются в дополнительном коде. 
                Рассмотрим 4-битное знаковое представление (диапазон от $-8$ до $7$). 
                Сложение максимального положительного числа $7$ ($0111$) и $1$ ($0001$) даёт:
                \[
                \begin{array}{cccc}
                0 & 1 & 1 & 1 \\
                + &   &   &   \\
                0 & 0 & 0 & 1 \\
                \hline
                1 & 0 & 0 & 0 \\
                \end{array}
                \]
                Результат $1000$ в дополнительном коде интерпретируется как $-8$. 
                В беззнаковой интерпретации $7+1=8$, но из-за переполнения мы получаем $-8$. 
                Таким образом, выход за пределы положительного диапазона отображается на отрицательные числа. 
                Если рассматривать бесконечность как число, большее максимального представимого, 
                то в замкнутой арифметике она ведёт себя как $-1$ в следующем смысле: 
                возьмём беззнаковое 4-битное представление (диапазон $0\ldots15$). Число $15$ ($1111$) — это $-1$ в дополнительном коде. 
                При сложении $15+1$ получаем:
                \[
                \begin{array}{|c|c|c|c|}
                \hline
                1 & 1 & 1 & 1 \\
                \hline
                \end{array}
                +
                \begin{array}{|c|c|c|c|}
                \hline
                0 & 0 & 0 & 1 \\
                \hline
                \end{array}
                =
                \begin{array}{|c|c|c|c|}
                \hline
                0 & 0 & 0 & 0 \\
                \hline
                \end{array}
                \quad\text{(перенос }1\text{ игнорируется)}.
                \]
                То есть $(-1) + 1 = 0$ по модулю $16$. 
                С точки зрения физики, ноль может быть суммой противоположных напряжений (например, +5В и -5В), 
                что соответствует аддитивности.
        \section{Модулярная арифметика}
                Современный компьютер использует сложение как базовую операцию (схема сумматора). 
                В модулярной арифметике основной операцией является умножение по модулю, 
                что упрощает реализацию многих алгоритмов (например, RSA, хеш-функции). 
                Переход к модулярному представлению позволяет выполнять умножение без переносов между разрядами.

                \textbf{Модель:} кольцо вычетов $\mathbb{Z}/n\mathbb{Z}$. 
                Операции выполняются по модулю $n$, причём умножение может быть быстрее при использовании специальных модулей 
                (например, чисел Мерсенна).
        \section{Простые числа и число 2}
                    Основная теорема арифметики утверждает, что любое натуральное число $n > 1$ можно разложить на простые множители (факторизовать) единственным способом, с точностью до порядка следования сомножителей.
                    \textbf{Единственное чётное простое.} Если бы существовало другое чётное простое $p>2$, то оно делилось бы на 2, а значит, имело бы делитель, отличный от 1 и себя, что противоречит определению. Это свойство делает двойку исключительной: все остальные простые числа нечётны.
                    \textbf{Основа двоичной системы счисления.} Современные компьютеры оперируют битами — состояниями 0 и 1. Двойка определяет алфавит цифр двоичной системы, которая лежит в основе всей цифровой техники и теории информации.
        \section{Дуплекс}
                Дуплекс — принцип двойственности: существует взаимно-однозначное соответствие между объектом и его представлением (информационной моделью). 
                Человек, видя объект, может создать его описание (словесное, графическое), а по описанию — восстановить объект в воображении или реальности.

                В информатике это соответствует сериализации и десериализации данных: объект $\to$ битовый поток $\to$ объект. 
                Также дуплекс проявляется в интерфейсах: кнопка (физическая) и её образ на экране (информационный).

                \textbf{Модель:} биекция $\phi: \mathcal{O} \to \mathcal{R}$, где $\mathcal{O}$ — множество объектов реальности, 
                $\mathcal{R}$ — множество их представлений (информационных моделей). 
                Обратное отображение $\phi^{-1}$ восстанавливает объект по модели.
\end{document}